\section{System architecture}
\label{sec:architecture}

This section describes the surface tactics, the side-channel architecture,
and the two-phase elaboration flow that lets users write cyclic proofs
interactively while producing a kernel-checked structural-induction theorem
as the user-facing declaration.

\subsection{Surface tactics}

CyclicTactic provides three primitive tactics, plus two command-level
forms.

\paragraph{\code{cyclic <label>}.}
Registers the current goal as a named \emph{companion}. Subsequent
back-edges within the same theorem body refer to this companion by label.
Internally, \code{cyclic} pushes a \code{.companion lbl seq} event into the
recorder (described below) and, for the first \code{cyclic} call in a body,
captures the goal head constant for use as \code{defaultSimpPred} when the
emitter needs to unfold the goal predicate.

\paragraph{\code{cyc\_cases x with | <pat> => \dots}.}
Case-splits on \code{x} while recording arm boundaries. Each arm's
syntactic range is captured so the recorder can later attribute back-edges
to the correct arm. The tactic delegates the actual case-split to Lean's
built-in \code{cases}, so InfoView goals appear as usual; the only
additional work is the event push.

\paragraph{\code{back <label> [\{σ\}]}.}
Closes the current goal via a recursive call to the enclosing theorem. The
\code{σ} substitution specifies how each binder of the enclosing theorem
maps to the call's arguments. \code{back R \{n := n', m := m\}} desugars to
\code{exact <name> n' m}, which Lean's standard structural-recursion check
accepts when \code{σ} produces strictly smaller arguments. The
back-edge's source position, σ, and surrounding source text are pushed as a
\code{.back} event.

\paragraph{Command-level forms.}
\code{cyclic\_thm name (binders) : type by tactics} packages the three
primitives into a command that produces a single user-facing declaration:
the kernel-checked structural-induction theorem produced by §6 emission, or
the fallback recursive form if §6 emission fails. \code{cyclic\_mutual
\dots end\_mutual} extends the same flow to mutually-recursive blocks
(§\ref{sec:impl-mutual}).

\subsection{Side-channel recording}

The cyclic structure of a proof---which case-splits, which back-edges, which
σs---is not visible to Lean's standard tactic machinery. Rather than
modifying Lean to track this information, CyclicTactic uses a
\emph{side-channel} architecture: each of the three primitive tactics pushes
an event into a global \code{CyclicState} \code{IO.Ref}; events are then
post-processed into a \code{ProofTree} value (a strongly-typed cyclic-proof
representation defined in \code{CyclicTactic/ProofTree.lean}).

The event-recording approach has three benefits over modifying Lean:
(1) the recorder is independent of Lean's elaborator internals, so it
survives Lean version upgrades; (2) the user gets ordinary InfoView goals
during proof writing, with no special editor support needed; (3) the
recorded structure is a pure-data Lean term, which we can validate, analyse,
and transform offline.

\begin{lstlisting}
inductive CyclicEvent where
  | companion (label : String) (sequent : Sequent)
  | caseSplitStart (var : String) (sequent : Sequent) (arms : List ArmInfo)
  | caseSplitEnd
  | back (ancestor : String) (sequent : Sequent) (σ : Subst)
         (pos : Nat) (sourceText : String)
\end{lstlisting}

The \code{ArmInfo} entries carry syntactic positions for each arm. The
\code{Build.eventsToTree} function (in \code{CyclicTactic/Build.lean})
walks the linear event stream: each \code{caseSplitStart}/\code{caseSplitEnd}
scope is converted to a \code{ProofTree.caseSplit}; back-edges within an arm
are mapped to that arm by source-position attribution.

\subsection{Two-phase elaboration}

The most subtle architectural choice is that \code{cyclic\_thm} elaborates
the proof \emph{twice} in succession: first as a recursive
\code{def}---which gives InfoView goals during interactive
writing---then, if §6 emission succeeds, as a structural-induction
\code{theorem} that becomes the user-facing declaration. The two phases
never coexist in the environment.

\begin{lstlisting}
-- Phase A: snapshot env before any work
let envBefore ← Lean.getEnv

-- Phase B: elaborate recursive form. User's tactics fire here;
-- InfoView shows real Lean goals at every cursor position.
Lean.Elab.Command.elabCommand
  (← `(def $name $binders : $type := by $tacs))

-- Phase C: snapshot env-with-recursive as fallback
let envWithRecursive ← Lean.getEnv

-- Phase D: build ProofTree from the recorded events;
-- run SCT + induction-order + §5 + §6 validations.
...

-- Phase E: try canonical §6 emission
let script := Theorem6.Emit.translate ...
match parse(script) with
| .ok cmdStx =>
  Lean.setEnv envBefore       -- roll back
  try elabCommand cmdStx       -- canonical form
  if hasGoodValue then ()      -- §6 emission committed
  else Lean.setEnv envWithRecursive  -- restore fallback
| .error _ =>
  -- recursive form stays as the user-facing decl
\end{lstlisting}

The user-facing \code{<name>} is always exactly one declaration. When §6
emission succeeds it is the structural-induction theorem; when §6 emission
fails for any reason it is the recursive \code{def} produced by Phase~B.

This architecture is what enables the system to be both ergonomic and
paper-faithful: the user writes cyclic-style with real interactivity, but
the kernel-checked artifact is the structural-induction proof Theorem~6.1
unravels. To our knowledge no prior cyclic-proof system in a mainstream
proof assistant uses this two-phase pattern.

\subsection{Pipeline overview}

The end-to-end flow from \code{cyclic\_thm} source to kernel-checked
theorem proceeds through nine stages, listed below with their owning
modules. \code{CyclicTactic/PIPELINE.md} walks the same flow on a
multi-recursive example (\code{btPredT}).

\begin{enumerate}
  \item \textbf{Command parsing and Phase~A env snapshot}
    (\code{Tactic.lean::elabCyclicThm}).
  \item \textbf{Phase~B recursive-form elaboration} (Lean's standard
    \code{def} machinery).
  \item \textbf{Event recording} (three custom tactics push into
    \code{CyclicState}).
  \item \textbf{ProofTree construction} (\code{Build.lean::eventsToTree},
    via source-position attribution).
  \item \textbf{SCT validation} (\code{ProofTree.extractTraceSCGsLabeled}
    \(+\) \code{SizeChange.checkMultiSCT}).
  \item \textbf{Paper-faithful §5 annotation}
    (\code{PaperAnnotation.lean}: Definition~5.1 step \(+\) Lemma~5.9
    unfolding \(+\) Theorem~5.2 verification).
  \item \textbf{§6 augmentation} (\code{Theorem6.lean::computeAug}:
    \(\RelAnc\)/\(\Ineq\)/\(\Hyp\)/\(\cov\)/\(\sigma\) per node).
  \item \textbf{§6 emission} (\code{Theorem6.Emit.translate}).
  \item \textbf{Canonical replacement} (Phase~E rollback and re-elaboration).
\end{enumerate}
